\documentclass[10pt]{article}
\usepackage[a4paper,margin=25mm]{geometry}
\usepackage[T1]{fontenc}
\usepackage{lmodern}
\usepackage{amsmath,amssymb,booktabs,array}
\usepackage[colorlinks=true,urlcolor=blue,citecolor=blue]{hyperref}
\usepackage{listings}
\lstset{basicstyle=\ttfamily\footnotesize,breaklines=true,columns=fullflexible,keepspaces=true,frame=single}
\title{ChoreoFix: Enforcing Multi-Agent Coordination from Global Protocols}
\author{Hayato Ohhashi}
\date{20 November 2026\\Research draft, comparative evaluation incomplete}
\begin{document}
\maketitle
\begin{abstract}
Agents can choose how to work, but shared effects still need clear rules.
We study whether a global protocol can serve as the source of those rules
without a separate application state machine. ChoreoFix uses Hibana to project
Rust global protocols into local communication programs. Handwritten async
roles drive their endpoints before carrying out domain effects. A bounded WASI
shell lets agents use work files and scripts. Every file operation goes through
ChoreoFS. The shell cannot grant a task step by choosing a pathname or writing
an approval file.
By 19 November, 36 of 48 tasks had qualified at the direct shell boundary.
For these tasks, 43 audited live-agent runs link shell commands and Hibana
operation occurrences to 732 shared effects and their replies. The common
simulator recorded task completion in 39 of these runs. Each qualification
applies to the software version that produced its trace. The sample includes
replacement runs. A subsequent seven-agent publication on task 7H added
a 37th qualified task. Formal benchmark scoring has yet to begin.
\end{abstract}

\section{Problem and scope}
A group of agents is a distributed program. An agent may prepare a change,
review a result, or choose another tool. It must still wait for the evidence or
resource needed by a shared operation. Putting an order in a prompt does not
make the effect depend on that order.

A runtime can check the order. The question here is how to define and maintain
that check. We write coordination as global protocols, project the local
communication programs, and let affine endpoints track progress. We keep
application computation in ordinary Rust. We do not maintain an application
program counter, completion flag, phase enum, or resource-owner map alongside
the endpoints.

The name ChoreoFix uses ``fix'' to mean fixing the permitted coordination
relation. It does not claim counterexample-guided automatic protocol repair.
Our research question is:
\begin{quote}
Can agents choose local work methods while projected global protocols guard
shared effects, retries, and independent work without a second application
progress authority?
\end{quote}

Three kinds of evidence are needed. A task comparison must measure safety and
completion on the original 48 tasks. A boundary experiment must show that
several legal command sequences work and that forbidden effects stop before
execution. A maintenance study must measure changes to declarations, localside
code, and authorization state. The present draft contains implementation and
integration evidence. The full task comparison and maintenance study are open.

\section{Background}
TraceFix synthesizes coordination protocols, checks PlusCal/TLA+ models, and
uses counterexamples to repair them. Its runtime study compares four ways to
execute the work \cite{tracefix}. We retain those original conditions. In
particular, a mediator already drives coordination. Runtime enforcement alone
is therefore not our claimed contribution. We test a different implementation
boundary and account for its limits.

Hibana takes direct inspiration from affine multiparty session types and
MultiCrusty \cite{ampst}. A global protocol describes communication among
roles. Projection derives each role's communication obligations. Hibana keeps
the resulting local program as data. An endpoint consumes permitted operations
without changing its Rust type after each step. Projection does not generate
the participant function. The function is handwritten, and its communication
is checked during execution.

\section{Execution boundary}
\subsection{Global and localside}
The implementation has three visible parts. Each task's \texttt{global.rs}
defines its workflow and handoffs. Its \texttt{localside/} functions show the
corresponding endpoint operations. The \texttt{imp/} code carries bytes and
model or simulator RPCs. It does not decide the next task step.

Sequential composition represents dependencies. Parallel composition retains
independent work. Routes represent alternatives, declared resolvers use actual
domain results, and repetition represents another attempt. These constructors
are used where the task needs them. A terminal test is not turned into a retry
loop merely to use a route.

A successful local step authorizes only the next operation represented by that
endpoint. It does not prove the quality of the agent's computation. The domain
operation is called by trusted localside code after the endpoint admits the
corresponding step. A failed operation cannot silently select another branch.
Task recovery follows an explicit route. A transport or session fault ends the
run and is not treated as a successful domain retry.

The review task, 6E, has two developers, a reviewer, and a shared repository.
The implementation assigns WASI roles, feature owners, reviewer services, and
one repository owner within the same global. The reviewer is one model and
one guest. That guest directly owns a caller endpoint for each service and a
private-I/O endpoint. Its two service owners do not add model participants.
Each feature can proceed while the other waits for its reviewer. A revision
uses the actual review result at an explicit resolver point:
\begin{lstlisting}
let feature_a = g::seq(
    implement_a,
    g::route(accepted_a, revise_a).resolve::<DECIDE_A>(),
).roll();
\end{lstlisting}
The complete global defines implementation requests, code-ready messages,
review requests and replies, repository operations, and completion. It composes
both features with \texttt{g::par}. The accepted arm admits merge and the
revision arm returns feedback before another implementation. Each corresponding
localside writes its \texttt{send}, \texttt{recv}, and \texttt{offer} operations
directly. For example, the developer consumes actual feedback:
\begin{lstlisting}
let feedback = endpoint.offer().await?;
match feedback.label() {
    30 => {
        let value = feedback.recv::<Approved<'_>>().await?;
        // Return this review result, then serve the merge.
    }
    31 => {
        let value = feedback.recv::<Revise<'_>>().await?;
        // Return this review result, then serve another implementation.
    }
    _ => return Err(...),
}
\end{lstlisting}
This is an excerpt, not a complete function. The source contains the actual
serialization and endpoint replies. Labels identify projected alternatives.
No approved flag or stored next-operation enum selects the continuation.
The resolver input holds the original review result, not task progress.

\subsection{Resources and notifications}
The sole repository owner receives projected requests before applying the
original implementation or merge operation. It owns the service object bytes
and awaits each domain operation before serving another. No separate lock
owner table grants access. This relies on the trusted single materializing
owner. Endpoint conformance does not by itself prove arbitrary Rust side
effects mutually exclusive.

Implementation and code-ready receipt precede review. The actual result reaches
the developer through the review-feedback channel. Approval must precede merge.
The global also carries the result to the guest. The trace auditor checks these
channels, effect intervals, and both completed merges separately from the
runtime's completion message.

The 16E CI/CD task connects its three WASI roles directly to the developer,
tester, deployer, and two resource roles in one global. Resource acquire, grant,
return, and acknowledgement are messages. Build completion precedes testing,
test completion precedes deployment, and deployment completion notifies both
other participants. Its original simulator has no injected domain failure.
A retry loop is therefore not added to this task.

Earlier codebase and writing prototypes connect task and resource sessions in
Rust brokers. Even when both sides use endpoints, that connection is outside
a single task global. Those prototypes and their tests are retained as component
evidence. They do not meet the target architecture and are not accepted task
ports. Porting must preserve disjoint work, original resource counts, and
terminal outcomes while replacing these connections. A serialized substitute
is not sufficient for tasks that allow disjoint operations.

\subsection{WASI and ChoreoFS}
The model sees one shell tool. Commands enter a real Wasm guest under
\texttt{hibana-wasip1-runtime}. Its pending WASI import is sent through a
projected ABI request/reply protocol. A matching typed completion resumes the
guest. ChoreoFS supplies path, object, fd, and I/O facts to the answering role.
Work files, scripts, terminal objects, and tool request/reply objects all use
this path. There is no guest-side substitute filesystem or host filesystem
fallback.

The current guest is the independent Rust \texttt{std} program
\texttt{hibana-wash}. It has no external Rust dependencies. It supports quoted
arguments, variables, conditions, loops, pipelines, redirection, calculations,
and saved scripts. Eight private files have 8,192-byte capacities. A script has
a 4,096-byte UTF-8 limit. Each guest starts with 64 KiB committed memory and can
request growth through the projected protocol up to 1 MiB of supplied backing.
These are explicit experimental limits, not an unrestricted POSIX environment.

The model receives object paths and the original domain schemas. It chooses
its commands. Writing JSON arguments to a domain object and closing it requests
that operation. The actual path is:
\begin{quote}
WASI endpoint $\rightarrow$ projected request $\rightarrow$ owner endpoint
$\rightarrow$ ChoreoFS operation $\rightarrow$ projected reply.
\end{quote}
The task global includes those edges. There is no unprojected request queue
connecting an ABI graph to a separate task graph. Private work uses a parallel
projected I/O fragment and cannot supply a missing task permission.

Object binding and argument validation require the target to match the chosen
service. A rejected argument follows an explicit resolver branch and can be
corrected without restarting the endpoint. ChoreoFS handles object operations,
not task authority. An open fd or an approval file cannot grant a merge.
The older limited-command guest and generic \texttt{/tools/call} bridge remain
only as historical component evidence.

\subsection{Trusted code and claim boundary}
The boundary assumes the supplied tool interface is the only path to the
protected effects. Model roles cannot call native tools or choose another
participant's identity. In 6E, 16E, and the 7H implementation, each simulator
call consumes a typed Hibana receive. Its type determines the participant and
operation, and its bytes determine the arguments. In 10M a compilation consumes
two requests, one from each required resource endpoint. The compiler compares
both with the submitted build arguments before making one simulator call.
Some ports still select an operation and supply its arguments to a privileged
simulator connection. In the selected 8E runs, the backend read the submitted
file contents and supplied the operation name at the call site. Those receivers
remain part of the trusted boundary. We found this interface in the archived
receiver sources for 3E, 7E, 8E, 13E/M/H and 14E/M/H. Together these nine
tasks account for 12 of the 43 selected conditions and 178 of their 732
simulator calls. Their successful run audits establish correspondence for the
recorded calls. The receiver API in those builds still permits a host caller
to supply the tool and arguments.
In a subsequent 8E build, all three simulator calls consumed typed receives.
Two additional three-agent trials passed the protocol and effect audits.
The original simulator scored the depth-zero trial complete and returned the
injected terminal test failure at depth two. The 43-condition analysis uses the
original builds.
All trials trust the runtime, codecs, transport, executor and simulator integration. The simulator
independently records actual domain effects and resource violations.

Endpoint conformance is not a proof of arbitrary program termination, review
quality, or artifact correctness. Existing Hibana and WASI proofs cover stated
models and assumptions. We do not claim a refinement proof for every line of
this application or a whole-system distributed deadlock theorem.

\section{Evaluation design}
\subsection{Original task comparison}
The selected matrix has 48 tasks, three failure depths, three repetitions,
and five conditions, for 2,160 scored episodes. The conditions are the paper's
Topology-monitored, Mediator-enforced, Prompt-only, Chat-only, and Hibana.
All participants use \texttt{gpt-6-luna} with xhigh reasoning. The episode limit
in the current matrix is 1,800 seconds and each participant has a 200-tool-call budget.

The four TraceFix runtimes remain at public revision
\href{https://github.com/Sensing-And-Reasoning/TraceFix/tree/87b01a2c0ae53c5045c9ce3d43cd0235b1831275}{\texttt{87b01a2}}. Their coordination logic is
unchanged. Benchmark descriptions and tool metadata at recorded revision
\href{https://github.com/Sensing-And-Reasoning/TraceFix/tree/36b4bd108fcd74f7efb17613aebbe3fdf55499a9}{\texttt{36b4bd1}} are byte-identical to that
release. A separate condition uses the unchanged latest monitoring runtime at
\texttt{36b4bd1}, adding 432 paired episodes across the same tasks, failure
depths, and repetitions. The total plan has 2,592 episodes. Its existing state,
monitor, distributed-coordination, and corrective-guidance tests passed in 94
cases. This is source qualification, not a measured agent comparison.
Our model and budgets differ from the original paper, so its percentages are not
matched controls for our runs.

Every port must retain the original roles, resources, dependencies, failure
paths, and required independence. All 825 checklist items remain in the
manifest. A passing compilation or generated TLC model is not by itself a
fidelity verdict. Task prototypes exist, but none has passed the
full independent benchmark acceptance gate. No scored episode has started.

\subsection{Which contention orders survive projection?}
An exclusive resource can be safe under several arbitration policies. In 9E the
upstream paper fixture orders the acquisition of the two forks for each diner.
We instrumented its unchanged TLA+ transition relation with the sequence of
meal events. TLC explored 152 distinct states; every permutation of A, B and C
reached a state in which all three had finished. The earlier ChoreoFix 9E global prescribed A, B, C. We changed the global to allow a table participant to select a ready meal request. Each diner still needs both fork endpoints before the effect. On one compiled projection, scripted Wasm guests exercised all six meal orders. Each run passed the EPF and effect audits, returned all six fork loans and completed in the shared simulator with zero violations. The first diner could finish before a second diner started to think. These runs used no model calls. They establish reachability of each meal order. The upstream fixture acquires the forks individually, so equality of all acquisition traces and fairness under sustained contention remain open.

In a separate unscored depth-zero trial on the same executable, three live
agents chose their own shell commands. The simulator recorded meals in B,A,C
order and scored the task complete with zero violations. The independent audit
checked all six fork returns and 1,578 EPF records from this execution.
Two TLC builds gave the same terminal orders on the upstream translated source. Their exact hashes and revisions are recorded in the reproduction. Qualification with the historical compiler pin remains open.

We then made the thinking operation repeatable in the 9E global and observed
both thoughts and meals in the upstream fixture. TLC exhausted 1,142 states
with this observer. We replayed all 90 distinct terminal action sequences via
the Wasm shells. Each run passed EPF and the request/effect audit. In an
additional run A thought three times, B and C ate between those thoughts, and
A ate last. The upstream fixture visits each think instruction once. On the
alphabet of participant-labelled thoughts and meals, the revised protocol
therefore admits a strict superset of the upstream fixture's terminal traces.
This statement concerns the three-participant task and hides fork-acquisition
steps. The extra behavior is a declared procedural difference for comparisons.
On the executable used for these replays, three live agents could choose
whether to repeat a thought. Each chose one. A began eating before B started
thinking, and meals ran in A,C,B order. All three finished, the simulator
recorded zero violations, and the audit checked 1,574 EPF observations and all
six fork returns. This unscored trial took 24.23 seconds.

The first five- and seven-participant versions with repeatable thinking
kept the earlier priority at contested forks. In scripted Wasm trials each diner recovered from an
invalid request, thought twice and ate once. Disjoint meals overlapped up to
two at a time in 9M and three in 9H. In separate live trials on that same
executable, all twelve participants chose to think once. Both trials completed
with zero violations and all ten and fourteen fork loans returned, respectively.
In 9H, A began eating before C and E first thought. On the executable used for both
live trials, all 90 upstream 9E action sequences and the additional sequence
with three thoughts by A again passed the EPF and effect audits and the common
task score.

The previous priority was observable at the shell boundary. In one archived depth-zero 9E
qualification, eat commands were dispatched in A,C,B order, while the simulator
started meals in A,B,C order. In 1E the first commit command came from B and
the first actual commit came from A. The archived 9M global allowed A and C,
or B and D, to eat concurrently and fixed precedence at each contested fork. Shell
dispatch time does not establish when every prerequisite for a meal or commit
was ready. We do not infer a scheduling penalty from these runs.

In the revised 9M and 9H globals, the first meal request received determines
which diner has first use of its two adjacent forks. Each remaining fork can
serve either neighbor first. We delayed the other guests until the first
guest had finished eating, rotating the first seat through all five or seven
positions. For each fixed first diner, the remaining diners ate in clockwise
order in one trial and counterclockwise order in another. All ten five-seat
and fourteen seven-seat Wasm trials passed EPF, the effect audit and the
common task score. For each first diner, the audited Endpoint receipts showed
both grant orders at every fork outside that diner's pair. Five live agents
on the same executable started eating in B,A,C,E,D order, with two eating at
once. Seven started in E,D,B,G,C,A,F order, with three eating at once. They
returned all ten and fourteen fork loans, respectively, with zero simulator
violations. These trials establish that any seat can go first. Fairness among
simultaneous arrivals remains unmeasured.

We removed the fixed priority in 1E. Each developer can submit its own
implementation independently. The repository receives the typed submission
and lends both required module endpoints to its sender. All three pairs of
developers share a module. Six scripted Wasm trials exercised every order,
with each developer finishing its local tests before the next began designing.
Each trial passed the common score and EPF audit and returned all six loans.
On the same executable, three live agents committed in C,B,A order. The
audit matched twelve effects and 4,957 protocol observations, and the common
simulator scored completion with zero violations. This unscored result shows
that the port preserves the choice of winner in these runs. It does not
measure fairness under simultaneous requests.

Can the agent choose again after a failure? In two five-guest Wasm runs of
10M, the same executable built modules A, B and C in that order. All three
needed another build. We held the failure injections and initial build order
fixed and varied only the scripted integrator's preference among eligible
modules. The simulator observed rebuilds A,A,B,C in one run and C,C,B,A in
the other. Both linked successfully with zero violations. The effect audit
matched all 15 operations in each run to their typed requests. EPF checked
48,589 and 46,070 observations, respectively. These are reachability tests
with scripted shells. Agent preferences were not measured.

In 4H, seven scripted Wasm guests exercised all 36 combinations of the three
services' schema migration and test orders on one executable. We also
injected a failure at each of the three test positions. In 27 of the 39
runs a test started before another service's review began. In 15 a test
started during another service's review. Every test named its own candidate
and a completed review of that candidate. The reviewer's shell received each
of the three repaired candidates. The runs ended before their second review.
The simulator recorded 510 effects and zero violations. All 39 task scores were incomplete because this
executable covered preparation and schema migration.

On a subsequent executable we injected a failure at each test position and
varied the review of the repair. For each position the reviewer approved the
first revision in one run and requested two further changes in another.
The six runs yielded twelve new review decisions. In each case the
reviewer received the service Endpoint, reviewed the fresh revision, and
returned the same decision to the reviewer and developer shells. In all
four runs where another service still had a test outstanding, one of those
tests completed between the failed test and the approval of its repair.
These runs ended before retesting. We assessed deployment and registration
orders separately in the common simulator. Given reviewed, tested candidates
and completed schema migrations, 720 of 2,250 graph/order pairs completed
without a violation across all 25 acyclic three-service dependency graphs.
Requiring every deployment to finish before the first registration would
exclude 432 of those 720 legal pairs. This count specifies behavior to
preserve when projecting the rest of the migration.

For each proposed mechanism comparison, we will first exercise alternate first
contenders, independent acquisitions, and an early request followed by a valid
retry in the same session. Both arms must expose the same admissible orders,
reply timing contract, and rejection and retry behavior. When they differ,
the native-system result can still be reported with its procedure; it cannot
isolate the effect of projection. The reproducible 9E observation and the
three archived dispatch sequences are described in \path{analysis/arbitration.md}.

\subsection{Outcomes and accounting}
Primary outcomes are full and partial completion and unauthorized committed
effects. We distinguish pre-effect rejection, domain failure, incomplete work,
timeout, protocol fault, and harness fault. We also record resource overlap,
latency, tool calls, and inclusive tokens. Cached inputs remain part of token
use. Reasoning tokens are not added twice.

The comparison must use paired task, failure-depth, and repetition identities.
Confidence intervals and paired contrasts will be computed only after scored
runs exist. Setup failures and unscored integration runs remain visible and
outside the scored denominators. A rejected effect is not a completed task.

\subsection{Separate shell experiment}
The shell changes the original tool interface. Its integration results are
therefore separate from the 2,160-run comparison. TraceFix keeps the original
interface. A causal comparison of shell coordination would require the same
guest, objects, prompts, and budgets in both arms.

The model receives the original task description, its assigned input data, and
the domain tool schemas. Some task prompts also describe role-specific choices
and recovery paths. These additions belong to the pilot condition. A paired
enforcement comparison must freeze the actual prompts, provide the same
information and exercise the same simulator version in both arms. Boundary tests
also exercise scripts and several legal file-operation sequences. A private
``approved'' file is deliberately written before an early shared operation.
The task endpoint still rejects the operation before the simulator receives
an execute request.

\section{Evidence to date}
\subsection{Historical source and core qualification}
The earlier integration snapshot used Hibana version \texttt{0.9.9} with an exact Git
patch to \texttt{c2001ec06944d7f3b299e16b5b3efb5b0ba0366a}. This is an unreleased
repair, not a new registry version. Both direct and WASI-runtime dependency
paths select that core. The WASI runtime is pinned to
\texttt{4ebdad1443cf49488eae924afe528f67f5a5953b} and wash to
\texttt{12e137aba2cea74f353c198a8d7c92afc977627c}.

Developing the 6E boundary exposed projection ambiguity at an elastic repetition
boundary and runtime problems involving parallel resolver visits and retained
send-selected choices. Minimal native regressions reproduce the affected cases.
The repair retains valid independent parallel repetitions. Local Rust tests
passed 883 cases with 12 explicit exclusions. The saved verification includes
Lean, Z3, focused Miri and Kani checks, embedded compilation, and Clippy.
These checks cover their recorded models and cases. They do not prove every
application trace safe. The core's separate full CI run 38734656111 also passed
both final-form and Kani at that exact commit.

A QUIC consumer selecting the same exact core passed 593 default and 598 HTTP/3
local tests. External run 38735739939 passed all 44 registered case/direction
cells at QUIC revision 44d291a in one run. Saved phase reports, verdicts,
expected cells, source-pair trees, and core/TLS pins were checked independently.
Reference self-tests were omitted. The retained evidence discloses runner and
capture-shutdown patches. This is integration evidence for the repair, not a
whole-system proof or a performance comparison.

\subsection{Early direct WASI model trials}
Table~\ref{tab:preflight} contains four unscored trials using the direct projected
boundary. Earlier limited-command and cross-graph bridge trials remain in their
separate ledgers. They are not evidence that the present architecture was
already complete. All participants in this table used the requested and returned
Luna model and xhigh effort. The original simulator and task inputs were kept.
\begin{table}[h]
\centering
\caption{Unscored direct WASI trials. No baseline contrast is available.}
\label{tab:preflight}
\begin{tabular}{lrrrrrr}
\toprule
Task & Depth & Roles & Shell calls & Domain calls & Seconds & Tokens \\
\midrule
16E & 0 & 3 & 11 & 3 & 27.127 & 220,215 \\
6E & 0 & 3 & 20 & 6 & 59.481 & 361,915 \\
6E & 1 & 3 & 23 & 8 & 45.387 & 412,438 \\
6E & 2 & 3 & 22 & 10 & 46.913 & 397,336 \\
\bottomrule
\end{tabular}
\end{table}

The 16E observation uses the pre-repair published core and is retained at its
own source checkpoint. All four required cross-role notifications and three
domain effects were audited. The three 6E observations use consumer 4a9937f and
core c2001ec. Each completed both merges with no recorded simulator violation.
The source snapshot was frozen before model calls and is identical across the
three review trials. Binaries, compiler details, prompts, raw model events, raw
carrier events, and inclusive usage were retained. Source hashes alone are not
used as evidence that a particular model turn finished.

For 6E the original simulator fails the first depth review calls across all
features. At depth 2 the model chose review order A, B, A, B. Each feature was
revised once. Deterministic boundary cases instead exercise A first and B first
with zero, one, or two revisions concentrated on the first feature. These cover
different observed schedules without requiring the model to follow a recipe.
The three model trials used 1,171,689 inclusive tokens in total. This is a cost
observation, not evidence of a token saving.

The read-only auditor replays the original simulator and checks every actual
close before its domain effect, nonoverlapping shared-resource intervals,
code-ready delivery, review-feedback contents, revised implementation after
rejection, merge after approval, and complete frame consumption. It also checks
raw model identity and turn completion, shell command correspondence, source
and binary hashes, and inclusive token arithmetic. Altered copies remove or
forge required requests and notifications, alter feedback, or overlap resource
intervals. They must fail even after the event-file hash is recomputed.

Later boundary tests found a caller-side limitation in that model checkpoint.
After a duplicate request to a completed review service, the shared caller
could still offer a legal branch of the other service. The application intended
a particular service; the public offer identifies an enabled endpoint branch.
This observation alone does not demonstrate that projection accepted two
indistinguishable public operations. All six added cases failed in the old application.
The revised global gives each service a separate caller role, directly driven
by the same original guest. All six repaired cases and simulator replays passed,
including chunked writes and duplicate requests followed by work on the other
feature. No application progress state or core change was needed. The earlier
model observations retain their original source and do not qualify this revised
source by themselves.

\subsection{Source-bound migration evidence}
The first 18 November verification batch covered 27 of 48 tasks across 21
families. Its mechanism suite contains 321 real Wasm conditions with
scripted shell inputs, including four retained CI/CD conditions whose
production-source continuity was checked. The separate live-agent suite
contains 28 required task--depth conditions across those same 27 tasks. The
accepted run for each condition has an independent audit. Across these runs,
the audits cover 119 model participants, 419 actual effects, and 312,880 records
from complete EPF captures. The recorded usage is 35,801,001 inclusive tokens.
The common simulator reported task completion in 27 of the 28 conditions.

The live ledger identifies the re-audited capture and the completed replacement
runs. It retains the interrupted audit and early guest exits from the original
attempts. The 28 condition results include replacements and do not estimate a
first-attempt success rate. The artifact records are
\path{output/verification/post-core-regression-20261117/v31/full-mechanism-regression.json},
\path{output/verification/post-core-regression-20261117/v31/live/verified-results.json},
and \path{output/verification/post-core-regression-20261117/v31/paper-inventory-recheck-v1.json}.
The last file checks the assigned task--depth pairs and links each accepted live
result to the hashes of its runtime, raw events, build receipt, and audit.

The separately qualified 13H port adds eleven Wasm conditions and three
live-agent conditions on one executable. The live runs involved five agents
each. All three passed the EPF and request/effect/reply audits. Two completed
the common-world task and one ran out of reagent stock. After final approval,
an additional blocking read had stalled an earlier live run. The final global
expresses the stock's finite retry paths. The same read returned a WASI error
in both a regression and a subsequent live run, allowing the guest to exit.
The completed conditions and their artifact hashes are recorded in
\path{output/experiments/lab-cascading-wasi-20261118/v42-finite-synthesis/verified-13H.json}.

Task 14E adds two real Wasm conditions and a four-agent live run. The
mechanism cases exercise both decisions and independent instrument effects.
In the live run, the lead read both actual reports, selected go, and the
chemist completed documentation. The common-world score and the EPF and
effect audits passed with 13 actual effects and zero violations.
The artifact is
\path{output/experiments/discovery-direct-20261118/v8-lead-evidence/verified-14E.json}.

Task 14M adds seven real Wasm conditions, including all four selective
reanalysis choices, two consecutive regulatory rejections, and retries after
failed instrument or documentation effects. A scientist's study file receives
reanalysis requested of that scientist or the actual approval. Unselected
scientists keep their existing evidence without resubmitting their reports.
Submitting a reanalysis choice requires no separate confirmation read.
Across the seven conditions the audits bind 135 effects to accepted requests
and replies, with 41,360 complete EPF observations, common-world completion,
and zero violations. The mechanism artifact is
\path{output/experiments/regulatory-discovery-direct-20261118/v5/verified-mechanism.json}.
The depth-two live run on that executable completed with five agents after
two regulatory rejections. The lead chose to reuse the current reports in both
revisions, performed a new structural confirmation, and submitted a new
decision for each review. Final approval and documentation passed common
scoring. The audit links 32 effects and 13,597 complete EPF records to the
same execution, with zero violations. The run took 329.59 seconds and used
2,666,036 inclusive tokens. The combined evidence is
\path{output/experiments/regulatory-discovery-direct-20261118/v5/verified-14M.json}.

Task 14H covers an entire material study with seven separate wash guests and
three consumable sample units. Three actual Wasm conditions completed direct
continuation, continuation after a two-instrument review, and documentation of
an evidence-backed abort after stock exhaustion. These executions produced
48 actual effects and 20,583 complete EPF records. The depth-two live run used
the same runtime binary and completed two rounds. Seven agents chose their
own retries and concluded with a documented abort once the evidence could no
longer support continuation. The common-world score and the EPF and effect
audits passed for 34 effects and 17,431 records with zero violations. That run
took 320.36 seconds and used 3,130,977 inclusive tokens. The same captured
record supported re-audit after EPF learned to construct a candidate with
unequal nested visit counts.
\path{output/experiments/multi-round-discovery-direct-20261118/v5/verified-14H.json}.

Task 3M uses five wash guests for three authors, a fact checker, and an
editor. Chapter Endpoints transfer during review and revision. The document
owner checks each verdict against the current draft before delivering it
to the author. Supplemental research remains independent until guest exit.
Two scripted Wasm conditions account for 40 effects and 15,913 complete
EPF records. In the depth-two live run, the five agents published the paper
after two factual and two editorial rejections. The same binary produced
26 effects, 22,434 complete EPF records, and two overlapping effects on
distinct physical resources. The common score and the audits passed with
zero violations in 174.78 seconds and 2,251,389 inclusive tokens. The
source-bound record is
\path{output/experiments/research-paper-direct-20261118/v5/verified-3M.json}.

Task 7M exercises an effect that needs two resources. A document Endpoint
receives a request to insert a particular figure in a chapter, forwards the
request to the figure store, and receives its confirmation. The receiver
consumes both incoming messages and compares their arguments before calling
the simulator. Six scripted runs with five Wasm guests reached publication
on the same executable. The three initial drafts completed in all six possible
orders. Five of the twelve insertions ran while at least one other chapter
still lacked a completed draft. Across 118 effects and 34,845 EPF records,
all 12 figure insertions had two
distinct committed receives before the corresponding simulator call. Changing
the figure ID in a saved store request caused the audit to reject that record.
The journal rows, effect IDs and build hashes are in
\path{analysis/7m-figure-receipts.json}.

A separate depth-one 7M trial with five live agents ended before publication.
The fact checker received the revised first chapter in its shell response.
It issued no further fact check and sent \texttt{exit 0}; the session then
closed after 17 effects, 1,744 seconds into its 1,800-second budget.
We tested that continuation on the same executable and Wasm guest. A
scripted fact checker tried rereading the consumed notice and result, received
permission errors, and submitted the revised pair of chapters for checking.
The document reached publication with 24 effects, zero scoring violations
and 7,884 records accepted by EPF. This case was selected after observing the
live failure. The two records are identified in
\path{analysis/7m-fact-recheck.json}.

In task 15E, one scripted wash guest used a delegated reference-wafer
Endpoint to mount the wafer, measure alignment, and run two spot checks in
either order. In three runs, the measurement started while the mount call was
still outstanding. In the recovery run, the first measurement failed. An
immediate spot check returned a shell error, and the guest measured the same
load again without remounting. Both final instrument results referred to the
original mount and the successful measurement. All three runs completed the
common-world score with zero violations. EPF checked the communication records;
we separately matched each instrument request and reply to the scored effect.
The evidence and timing conditions are in \path{analysis/15e-wafer-alignment.json}.
A subsequent six-run fixture added three wash guests to receive the engineer's
lot disposition and the two inspection receipt identifiers. All six send
orders completed. In three runs all writes finished before the first read,
and the recipients read in reverse order. Across the six runs EPF checked
11,691 records, and the common simulator recorded zero violations.
The other five roles' upstream work remained fixture-driven; these trials
made no model calls. Per-run receipts, delivery order and hashes are in
\path{analysis/15e-wafer-disposition.json}.

\subsection{Live-agent evidence across 36 tasks}
The 19 November inventory includes one audited execution for each of the
43 qualified task--depth conditions across 29 families. We matched each raw
event stream to its saved audit, executable, source manifest and participant
reports. For each of the 2,466 shell calls, the model server acknowledged the
same tool arguments and reply contents as the driver recorded for that call.
The runs contain 197 participant sessions, 563,288 EPF observations,
and 67,938,442 inclusive tokens. The four negative common-world scores were
8E at depth two, 13H at depth two, and 16M at depths one and two. All four
passed the protocol and effect audits. The selected executions, including
replacements, are individually identified in
\path{output/verification/paper-36-ports-20261119.CogEVn/verified-conditions.json}.
The archived manifests contain two distinct sets of Hibana core sources.
Both predate changes to receive admission and offer handling used in the
subsequent 7H trial. The 43 runs establish results for their recorded
builds. The per-run source comparison is archived in
\path{output/verification/paper-36-ports-20261119.CogEVn/core-source-comparison.json}.

Across these 43 runs, we measured all 732 simulator calls on the runner's
monotonic clock. Calls from different participants overlapped in 27 runs
covering 26 distinct tasks. There were 318 overlapping pairs. The peak number
of participants with a call in progress was one in 16 runs, two in 17,
three in six, and four to six in the remaining four. The seven-developer
1H run at depth two reached six. Each interval extends from the start of
a domain request through its simulator callback, including any simulated
delay. Two intervals count as overlapping only if their intersection has
positive duration. Per-run bounds and request IDs are archived in
\path{output/verification/paper-36-ports-20261119.CogEVn/effect-overlap.json}.

All 732 simulator callbacks returned a Boolean outcome. In 20 of the
43 selected conditions, the simulator returned a negative outcome at least
once, for a total of 112 negative outcomes. For each we located the first
later call by the same participant to the same tool whose start followed
the earlier call's end. There were 102 such follow-ups. In 76 the JSON
arguments were equal; 42 of those calls succeeded and 34 returned another
negative result. Among the 26 with different arguments, 12 succeeded.
Ten negative outcomes had no such follow-up. A negative outcome includes,
for example, a review requesting changes. Matching the tool alone can
also match work on a different artifact. These counts describe actions
in the selected runs and do not estimate a recovery probability.

The five-agent CI/CD task (16M) illustrates what the two outcomes measure.
With no injected failure, the pipeline passed the production smoke test in
100.11 seconds (14 effects). At depth one, staging deployment failed after
26 effects in 273.40 seconds. At depth two, the second QA check failed after
30 effects in 314.97 seconds. Both failures produced a negative business
score. The audit found zero protocol or effect violations in all three runs.
\path{output/verification/cicd-16m-20261119/verified-live-depths-0-1-2.json}.

None of the 70 simulator calls in those three runs overlapped a call from
another participant.
\path{output/verification/paper-36-ports-20261119.CogEVn/16m-effect-intervals.json}.

\subsection{Limits of the integration evidence}
Task 7H exposes a different outcome. In two seven-agent trials outside the
qualified set, the reviewer issued \texttt{exit 0} while the manuscript
still needed work. At depth one, its preceding read of the editorial file
returned a permission error. At depth two, it had received an explicit
request for another revision of Chapter 4. In both traces, the response
delivered to the model matched the guest's terminal output. The simulator
recorded 14 and 20 effects respectively, zero violations, and an unfinished
document. A shell that rejects unavailable operations still lets an agent
choose to leave early.
\path{output/verification/7h-schema-guide-20261119.FTJvex/reviewer-exit-depths-1-2.json}.

At depth zero, the editor read the material before the manuscript was ready.
The model server cancelled the outstanding tool call after 471.984 seconds.
The guest subsequently returned 112,435 bytes. The driver sent that output
26.96 seconds after the server had closed the call with an empty result.
The successful read in the guest trace therefore did not establish that the
editor received the manuscript. Its last required review finished only
1.82 seconds before wash returned the reviewed material. Four participants
reached the 1,800-second group deadline with an outstanding shell exit. For
each of the four, a shell call was outstanding during 86--92\% of its turn.
The remaining interval includes model generation and server scheduling.
The source record identifies the command and the server call separately and
matches the guest bytes to the late reply.
\path{output/verification/7h-schema-guide-20261119.FTJvex/material-delivery-depth-0.json}.

In a subsequent depth-zero trial with seven agents, the authors and reviewers
completed three rounds of substantive revision on Chapter~4. After acceptance,
the fact checker checked Chapter~3. At 390 seconds the reviewer was already
waiting on its editorial-outcome file. At 398 seconds the transport recorded
the Chapter~3 evidence going to the reviewer's material file (frame 20408).
The guest returned no Chapter~3 manuscript to the reviewer. The trial reached
its 1,800-second deadline with 24 effects and two completed participants.
The common simulator saved an incomplete score. The end-to-end audit rejected
the trial. This run used the ordinary world and delays, with no injected
business failure.
\path{output/verification/7h-terminal-outcome-20261120/verified-results.json}.

In a later seven-agent 7H trial, the reviewer requested a major Part~I
revision, inspected the new fact check and approved the manuscript. The
editor published it. The common simulator recorded 36 effects and zero
violations. Two of the editor's shell calls had been interrupted. The model
connection acknowledged the recovered output of the first. The second
output remained uncollected when the editor's turn ended, and the runner
stopped before the final capture checkpoint. The trial therefore has a
positive business score and an unsuccessful end-to-end audit.
\path{analysis/7h-review-inbox.json}.

In a scripted 7H trial, the reviewer read its inbox twice before submitting
the next assessment. On the earlier protocol, the second blocking read
received no reply and the trial timed out after 36 simulator effects.
We revised the global so that each Part~II chapter's fact check and review
use one manuscript endpoint, returned to the drafting participant on approval.
With the same shell commands and review choices, both the fact checker
and the reviewer received a permission error on their premature reads
and continued to submission. The trial finished with 39 effects.
Across all 18 scripted trials on this binary, the common simulator marked
every manuscript complete. The audits matched 565 effects and checked
272,879 EPF records with zero violations. These trials used no model calls.

Seven live agents then used the same binary. After two revisions, the
Chapter~3 writer tried to replace the chapter text following another
negative review. The shell rejected the write. The writer exited, and the
run ended after 22 effects without a complete manuscript or a passing
audit. Replaying those 22 effects in the common simulator and supplying the
same replacement text produced a successful write. The global permitted
appending revision feedback at this point, but omitted the writer's choice
to replace the text. This trial remains outside the qualified set.
\path{analysis/7h-review-inbox.json}.

On the revised global, each Part~II author can respond to a major revision
by appending feedback or replacing the chapter through its original file.
In a new Wasm trial, the authors of Chapters~3 and~4 each appended once,
then replaced the entire text after a second rejection. Both replacements
received fresh fact checks and reviews. The Part~I author read the exact
replacement texts, and the common simulator accepted the final manuscript.
All 19 scripted trials on this binary completed, with 604 effects and
289,212 EPF records checked. The trials used no model calls.

Seven live agents subsequently used the same binary at depth zero. They
completed the manuscript in 842.61 seconds and used 8,352,961 inclusive
tokens. Three negative reviews led to further edits. The authors used the
append operation for their Part~II revisions. All seven participants
completed, and the audit matched 33 effects and replies to the same trace
of 147,787 EPF records. The common simulator recorded publication and zero
violations. This additional trial is outside the November~19 cohort.
\path{analysis/7h-review-inbox.json}.

These are integration observations from development, not random samples or
paired baseline trials. They cannot estimate a comparative completion rate,
latency advantage, or maintenance saving. Earlier failed and interrupted
attempts are retained in the local experiment archive. The public source
repository contains summaries and reproduction commands. It does not include
the raw journals for independent reanalysis.

The 43 selected ChoreoFix runs used 67,938,442 inclusive tokens in total (median 1,027,748 per run). The inventory selects qualified conditions; it does not estimate the token budget or success rate of another arm.

The shell changes the domain-call interface. In 6E a developer's implementation
reply waits for review and includes its actual result. Canonical object targets
are checked before execution. The unchanged baseline retains its own argument
matching and original tool replies. These differences must be controlled in a
causal shell comparison or reported as a separate interface experiment.
For the controlled comparison, both arms must receive the same task and
role information. We will freeze the prompts and simulator revision and run
a smoke trial for each included arm and task before scoring. Current live
pilots with task-specific shell guidance cannot isolate the effect of endpoint
enforcement.

The original tools simulate work and do not consume an agent's source patch.
They can measure coordination, not whether one reasoning strategy produces a
better software repair. Private shell computation is flexible within the
stated profile. It does not make this a general operating-system sandbox or a
proof of unrestricted agent freedom. No complete 48-task fidelity acceptance,
scored comparison, or maintenance measurement is reported here.

\section{Maintenance study and remaining work}
We specify three paired changes before collecting maintenance results: add a
second candidate-bound approval in 6E; add a replacement-and-retest path after
a failed deployment check in 16E; and allow all first diners from a matched
fixed-order starting procedure in 9E. The 9E upstream fixture already permits
all six orders. An unmatched starting pair is recorded as such. We will freeze
the starting revisions, editing budgets, editor assignment and the same
observable acceptance cases for each pair before editing.

The primary outcome per ticket is the number of independently maintained
authorization decisions that must change. Two manual rules that can disagree
about a request count separately even in the same file. Two reviewers will
label edits independently and reconcile disagreements. Secondary outcomes are
edited files, engine versus application changes, elapsed time, and remaining
acceptance failures. The change cost includes edits to task-specific auditors,
scripted guests, and regression tests. We report those edits separately from
the global, localsides, common simulator, and engine. Generated artifacts are
accounted for separately.
Unfinished tickets retain their time and failing observations. Ticket-level
results and missing pairs will be reported individually. No maintenance
advantage has yet been measured.

The main blockers to a completed empirical paper are completion of all task
ports and independent fidelity acceptance of all 48, a runnable and accepted
baseline corpus, the 2,160 main and 432 supplemental scored episodes, analysis
with uncertainty, and the
maintenance measurements. The shell profile also limits the kinds of work an
agent can carry out. The original simulator checks coordination outcomes and
does not establish the quality of arbitrary generated source code.

\section{Conclusion}
ChoreoFix provides an executable way to keep workflow progress in projected
endpoints while agents choose commands in a ChoreoFS-backed WASI shell.
The present evidence covers resource handoffs, explicit test retries,
independent work, and rejection before a forbidden task effect. It supports
further evaluation. It does not yet support a comparative performance claim or
a submission-ready empirical conclusion.

\begin{thebibliography}{9}
\bibitem{tracefix} Shuren Xia, Qiwei Li, Taqiya Ehsan, and Jorge Ortiz.
\emph{TraceFix: Repairing Agent Coordination Protocols with TLA+ Counterexamples}.
2026. \url{https://arxiv.org/abs/2605.07935}.
\bibitem{ampst} Nicolas Lagaillardie, Rumyana Neykova, and Nobuko Yoshida.
\emph{Stay Safe under Panic: Affine Rust Programming with Multiparty Session Types}.
ECOOP 2022. \url{https://arxiv.org/abs/2204.13464}.
\end{thebibliography}
\end{document}
